\chapter{CNN}
\section{CNN简介}
本章的主题是卷积神经网络（Convolutional Neural Network，CNN），
之前介绍的神经网络中，相邻层的所有神经元之间都有连接，这称为全连接网络（fully-connected Network）。
全连接网络主要由全连接层组成，其本质上是线性变换加偏置（$y = Wx + b$），所以也叫做 Affine 层。

\begin{figure}[H]
	\centering
	\includegraphics[width=0.9\textwidth, keepaspectratio]{3-1.pdf}
	\caption{全连接网络示例}
\end{figure}

相较全连接网络，CNN中新出现了卷积层（Convolution层）和池化层（Pooling层），
CNN 的层的连接顺序是“Conv - ReLU - Pooling”（Pooling层有时会被省略）。
还需要注意的是，在 CNN 中，靠近输出的层中使用了之前的“Affine - ReLU”组合。
此外，最后的输出层中使用了之前的“Affine - Softmax”组合，这些都是 CNN 中比较常见的结构。

\begin{figure}[H]
	\centering
	\includegraphics[width=0.9\textwidth, keepaspectratio]{3-2.pdf}
	\caption{卷积神经网络示例}
\end{figure}

\section{卷积层}
在介绍卷积层之前，我们先来看看全连接层（Affine层）在处理图像数据时的局限。
图像通常是高、长、rgb通道方向上的3维形状，其中包含了重要的空间信息——相邻像素往往具有相似的值，形状中也可能隐藏着值得提取的特征。
然而，向全连接层输入图像时，需要将3维数据拉平为1维数据，这些空间信息通通被“忽视”了。

为了解决这个问题，CNN 引入了卷积层。卷积层可以保持形状不变，当输入数据是图像时，卷积层会以3维数据的形式接收输入数据，并同样以3维数据的形式输出至下一层。
利用卷积层，有可能正确理解输入数据的形状的信息。卷积层的输入输出数据称为特征图（feature map），其中，输入数据称为输入特征图（input feature map），输出数据称为输出特征图（output feature map）。

卷积层进行的处理就是卷积运算，相当于图像处理中的“滤波器运算”。
具体来说，卷积运算会以一定间隔滑动滤波器（也称为卷积核）的窗口，对窗口内的元素进行乘积累加运算。

\begin{figure}[H]
	\centering
	% 第一行
	\begin{subfigure}[b]{0.4\textwidth}
		\centering
		\includegraphics[width=0.8\textwidth]{3-3a.pdf}
		\caption{step 1}
		\label{fig:a}
	\end{subfigure}
	\hfill
	\begin{subfigure}[b]{0.4\textwidth}
		\centering
		\includegraphics[width=0.8\textwidth]{3-3b.pdf}
		\caption{step 2}
		\label{fig:b}
	\end{subfigure}

	\vspace{0.5cm} % 行间距

	% 第二行
	\begin{subfigure}[b]{0.4\textwidth}
		\centering
		\includegraphics[width=0.8\textwidth]{3-3c.pdf}
		\caption{step 3}
		\label{fig:c}
	\end{subfigure}
	\hfill
	\begin{subfigure}[b]{0.4\textwidth}
		\centering
		\includegraphics[width=0.8\textwidth]{3-3d.pdf}
		\caption{step 4}
		\label{fig:d}
	\end{subfigure}

	\caption{卷积运算流程说明}
	\label{fig:grid}
\end{figure}

CNN中也存在偏置，不过偏置通常只有1个，这个值会被加到应用了滤波器的所有元素上。

\begin{figure}[H]
	\centering
	\includegraphics[width=0.8\textwidth, keepaspectratio]{3-4.pdf}
	\caption{偏置示例}
\end{figure}

在进行卷积层的处理之前，有时要向输入数据的周围填入固定的数据（比如0），这称为填充（padding），主要是为了调整输出的大小。
如图\ref{fig:padding}所示，通过应用幅度为1的填充，(4, 4)的输入数据变成了(6, 6)的形状，
然后，应用形状为(3, 3)的滤波器，生成了形状为(4, 4)的输出数据。

\begin{figure}[H]
	\centering
	\includegraphics[width=0.6\textwidth, keepaspectratio]{3-5.pdf}
	\caption{填充示例}
	\label{fig:padding}
\end{figure}

应用滤波器的位置间隔称为步幅（stride），图 \ref{fig:stride} 为一个步幅为2的卷积运算示例。

\begin{figure}[H]
	\centering
	\begin{subfigure}[b]{0.4\textwidth}
		\centering
		\includegraphics[width=0.8\textwidth]{3-6a.pdf}
		\caption{step 1}
	\end{subfigure}
	\hfill
	\begin{subfigure}[b]{0.4\textwidth}
		\centering
		\includegraphics[width=0.8\textwidth]{3-6b.pdf}
		\caption{step 2}
	\end{subfigure}

	\caption{步幅示例}
	\label{fig:stride}
\end{figure}

图像是3维数据，除了高、长方向之外，还需要处理通道方向，在3维数据的卷积运算中，输入数据和滤波器的通道数要设为相同的值。
最后数据输出是1张特征图，如果要数据输出在通道方向上也拥有维度，那么就需要用到多个滤波器，生成多个特征图。

\section{池化层}
经过卷积层之后，特征图的空间尺寸可能仍然很大，这会导致后续层的参数量和计算量过多。
池化层的作用就是缩小高、长方向上的空间尺寸，从而减少计算量，同时保留主要的特征信息。
下图是按步幅2进行2 × 2的Max池化时的处理顺序。
一般来说，池化的窗口大小会和步幅设定成相同的值。
除了Max池化之外，还有Average池化、随机池化等，在图像识别领域，主要使用Max池化。

\begin{figure}[H]
	\centering
	% 第一行
	\begin{subfigure}[b]{0.4\textwidth}
		\centering
		\includegraphics[width=0.8\textwidth]{3-7a.pdf}
		\caption{step 1}
	\end{subfigure}
	\hfill
	\begin{subfigure}[b]{0.4\textwidth}
		\centering
		\includegraphics[width=0.8\textwidth]{3-7b.pdf}
		\caption{step 2}
	\end{subfigure}

	\vspace{0.5cm} % 行间距

	% 第二行
	\begin{subfigure}[b]{0.4\textwidth}
		\centering
		\includegraphics[width=0.8\textwidth]{3-7c.pdf}
		\caption{step 3}
	\end{subfigure}
	\hfill
	\begin{subfigure}[b]{0.4\textwidth}
		\centering
		\includegraphics[width=0.8\textwidth]{3-7d.pdf}
		\caption{step 4}
	\end{subfigure}

	\caption{池化层流程说明}
\end{figure}

池化层有几个重要特征。首先，和卷积层不同，池化层没有要学习的参数。
其次，经过池化运算，输入数据和输出数据的通道数不会发生变化。
此外，池化层对微小的位置变化具有鲁棒性——即使输入数据发生微小偏差，池化仍会返回相同的结果。

\section{CNN的发展}
\subsection{AlexNet}
AlexNet 由 Alex Krizhevsky 等人于 2012 年提出，在 ILSVRC 这个大规模图像识别比赛中以显著优势夺冠，标志着深度学习在计算机视觉领域的崛起。
其主要创新包括：使用 ReLU 激活函数替代传统的 Sigmoid，缓解了梯度消失问题；引入 Dropout 随机失活来抑制过拟合；以及利用 GPU 进行并行训练加速。
网络结构上，AlexNet 包含 5 个卷积层和 3 个全连接层，整体较之前的网络更深。

\subsection{VGG}
VGG 由牛津大学视觉几何组（Visual Geometry Group）于 2014 年提出，其核心思想是通过堆叠多个小卷积核（$3 \times 3$）来替代单个大卷积核（如 $5 \times 5$、$7 \times 7$）。
这样做有两个好处：一方面，多个 $3 \times 3$ 卷积核的感受野与一个大卷积核相同，但参数量更少、计算更高效；另一方面，更多的非线性层增强了网络的表达能力。
VGG 网络（如 VGG16、VGG19）通过这种简洁统一的设计，将网络深度提升到了 16 层甚至 19 层，在 ILSVRC 中取得了优异成绩。

\subsection{GoogLeNet}
GoogLeNet 同样在 2014 年的 ILSVRC 中表现出色，但其设计思路与 VGG 不同，它的特征是，网络不仅有深度，而且有广度。
实现广度的是 Inception 模块：在同一层中并行使用多种不同尺寸的卷积核（$1 \times 1$、$3 \times 3$、$5 \times 5$）以及池化操作，然后将它们的输出拼接在一起。
这样网络可以同时捕捉不同尺度的特征，而无需人工决定使用哪种卷积核。
此外，GoogLeNet 大量使用 $1 \times 1$ 卷积核来进行降维，减少计算量和参数量。

\subsection{ResNet}
随着网络层数的增加，一个看似矛盾的问题出现了：更深的网络反而表现更差——训练误差和测试误差都上升了，这并非过拟合所致，而是网络难以被优化。
ResNet（残差网络）由何恺明等人于 2015 年提出，通过引入“小路”（skip connection）巧妙地解决了这个问题。
其核心思想是“残差学习”：不再直接学习目标映射 $H(x)$，而是学习残差 $F(x) = H(x) - x$，最终输出变为 $F(x) + x$。
这种结构使得梯度可以绕过中间层直接回传，有效缓解了深层网络中的梯度消失问题。
通过在VGG中引入这个设计，ResNet 成功训练了 152 层的超深网络，并在 ILSVRC 中大幅领先。
